\documentclass[10pt,a4paper,landscape]{article}

% --- Paquetes ---
\usepackage[utf8]{inputenc}
\usepackage[spanish,es-noshorthands]{babel}
\usepackage[margin=1.5cm]{geometry}
\usepackage{xcolor}
\usepackage{tikz}
\usetikzlibrary{positioning,calc,shapes.geometric}
\usepackage{enumitem}
\usepackage{fontawesome5}
\usepackage{tabularx}
\usepackage{booktabs}
\usepackage{multicol}
\usepackage{hyperref}

% --- Colores ---
\definecolor{accentblue}{HTML}{2563EB}
\definecolor{accentred}{HTML}{DC2626}
\definecolor{accentgreen}{HTML}{16A34A}
\definecolor{lightgray}{HTML}{F3F4F6}
\definecolor{darktext}{HTML}{1F2937}
\definecolor{medgray}{HTML}{6B7280}

% --- Sin numeración de páginas ---
\pagestyle{empty}

% --- Fuente ---
\renewcommand{\familydefault}{\sfdefault}

% --- Comandos de estilo ---
\newcommand{\sectiontitle}[2]{%
  \noindent\colorbox{accentblue}{\parbox{\dimexpr\linewidth-2\fboxsep}{%
    \centering\color{white}\bfseries\large #1\quad #2%
  }}%
  \vspace{6pt}%
}

\newcommand{\keystroke}[1]{%
  \tikz[baseline=(key.base)]{%
    \node[fill=lightgray, draw=medgray, rounded corners=3pt, 
          inner xsep=3pt, inner ysep=2pt, font=\bfseries\ttfamily\scriptsize] (key) {#1};%
  }%
}

\newcommand{\warningbox}[1]{%
  \noindent\colorbox{accentred!10}{\parbox{\dimexpr\linewidth-2\fboxsep}{%
    \color{accentred}\bfseries\faExclamationTriangle\quad #1%
  }}%
}

\newcommand{\tipbox}[1]{%
  \noindent\colorbox{accentgreen!10}{\parbox{\dimexpr\linewidth-2\fboxsep}{%
    \color{accentgreen!80!black}\bfseries\faCheckCircle\quad #1%
  }}%
}

\begin{document}

% --- Título ---
\begin{center}
  {\color{accentblue}\rule{\linewidth}{2pt}}\\[6pt]
  {\Huge\bfseries\color{darktext} \LaTeX: Overleaf + VS Code}\\[4pt]
  {\large\color{medgray} Tarjeta de referencia rápida}\\[4pt]
  {\color{accentblue}\rule{\linewidth}{2pt}}
\end{center}

\vspace{8pt}

\begin{multicols}{3}

% ============================================================
% COLUMNA 1: Flujo principal + Atajos
% ============================================================

\sectiontitle{\faRoute}{Flujo de trabajo}

\begin{center}
\begin{tikzpicture}[
  node distance=0.6cm,
  every node/.style={font=\small},
  box/.style={rectangle, rounded corners=5pt, draw=accentblue, 
              fill=accentblue!8, text width=4.2cm, align=center,
              minimum height=0.7cm, font=\small},
  decision/.style={diamond, draw=accentred, fill=accentred!8,
                   text width=2cm, align=center, font=\small\bfseries,
                   inner sep=1pt},
  arrow/.style={->, thick, >=stealth, color=medgray}
]
  \node[box] (work) {Trabajar en \textbf{Overleaf}};
  \node[decision, below=of work] (broke) {¿Se rompió?};
  \node[box, below left=0.5cm and -0.8cm of broke] (no) {Seguir en Overleaf};
  \node[box, below right=0.5cm and -0.8cm of broke] (yes) {Pasar a VS Code};
  \node[box, below=of yes] (fix) {Arreglar y resubir};
  
  \draw[arrow] (work) -- (broke);
  \draw[arrow] (broke) -- node[left, font=\scriptsize] {No} (no);
  \draw[arrow] (broke) -- node[right, font=\scriptsize] {Sí} (yes);
  \draw[arrow] (yes) -- (fix);
  \draw[arrow] (fix.south) .. controls +(0,-0.4) and +(0,-0.4) .. 
    node[below, font=\scriptsize] {Volver} +(-2.5,0) |- (work.west);
\end{tikzpicture}
\end{center}

\vspace{8pt}

\sectiontitle{\faKeyboard}{Atajos en VS Code}

\vspace{2pt}

\renewcommand{\arraystretch}{1.6}
\begin{tabularx}{\linewidth}{@{}lXX@{}}
  \toprule
  \textbf{Acción} & \textbf{Windows} & \textbf{Mac} \\
  \midrule
  \faHammer~Compilar & \keystroke{Ctrl+Alt+B} & \keystroke{Cmd+Opt+B} \\
  \faFilePdf~Ver PDF & \keystroke{Ctrl+Alt+V} & \keystroke{Cmd+Opt+V} \\
  \faSave~Guardar & \keystroke{Ctrl+S} & \keystroke{Cmd+S} \\
  \bottomrule
\end{tabularx}

\vspace{6pt}

\tipbox{Después del primer Build, Ctrl+S guarda Y compila automáticamente.}

% ============================================================
% COLUMNA 2: Pasos para rescate + Regla de oro
% ============================================================

\sectiontitle{\faLifeRing}{Rescate: Overleaf → VS Code}

\begin{enumerate}[leftmargin=*, itemsep=3pt, font=\bfseries\color{accentblue}]
  \item \textbf{Descargar} de Overleaf:\\
    {\small Menu → Download → \textbf{Source} (.zip)}
  \item \textbf{Descomprimir} el .zip:\\
    {\small Win: clic derecho → Extraer todo\\
    Mac: doble clic}
  \item \textbf{Abrir} en VS Code:\\
    {\small File → \textbf{Open Folder} → la carpeta}\\
    {\small\color{accentred} ¡La carpeta, NO el archivo!}
  \item \textbf{Compilar}: atajo de Build\\
    {\small Primera vez tarda 1--3 min (normal)}
  \item \textbf{Ver PDF}: atajo de Ver PDF
  \item \textbf{Arreglar} el problema y guardar
  \item \textbf{Subir} a Overleaf:\\
    {\small Upload → seleccionar archivos modificados → reemplazar}
\end{enumerate}

\vspace{6pt}

\warningbox{El PDF SOLO se ve dentro de VS Code.}

\vspace{4pt}

\noindent{\small\color{darktext} Si el PDF está abierto en \textbf{WPS Office}, \textbf{Adobe} u otro programa, VS Code no puede actualizarlo. Cierra el PDF en el otro programa primero.}

% ============================================================
% COLUMNA 3: Troubleshooting
% ============================================================

\sectiontitle{\faWrench}{Solución de problemas}

\vspace{2pt}

\renewcommand{\arraystretch}{1.4}
{\small
\begin{tabularx}{\linewidth}{@{}>{\raggedright\arraybackslash}p{2.8cm}X@{}}
  \toprule
  \textbf{Problema} & \textbf{Solución} \\
  \midrule
  PDF no se actualiza & 
    Cerrar PDF en WPS/Adobe. Solo verlo en VS Code \\
  \addlinespace
  Citas como~{[?]} & 
    1) Verificar \texttt{\textbackslash bibliography\{nombre\}} coincida con archivo .bib\newline
    2) Recompilar con atajo de Build \\
  \addlinespace
  Citas se rompieron de repente &
    Revisar el .bib: llaves \{ \} cerradas, comas entre campos \\
  \addlinespace
  Paquete no encontrado &
    MiKTeX lo instala solo. Si no: MiKTeX Console → Packages \\
  \addlinespace
  Primera compilación lenta &
    Normal. Esperar. Después será rápido \\
  \addlinespace
  Errores rojos abajo &
    Clic en el número → ver línea del error → corregir \\
  \bottomrule
\end{tabularx}
}

\vspace{10pt}

\sectiontitle{\faBookOpen}{Recuerda}

\begin{itemize}[leftmargin=*, itemsep=4pt]
  \item[\color{accentgreen}\faCheckCircle] \textbf{Overleaf} es tu herramienta principal
  \item[\color{accentblue}\faCheckCircle] \textbf{VS Code} es tu plan B cuando algo se rompe
  \item[\color{accentgreen}\faCheckCircle] \keystroke{Ctrl+S} / \keystroke{Cmd+S} = guardar + compilar
  \item[\color{accentred}\faTimesCircle] \textbf{Nunca} abrir el PDF fuera de VS Code
  \item[\color{accentgreen}\faCheckCircle] Si citas salen como {[?]} → recompilar
\end{itemize}

\end{multicols}

\end{document}
